\chapter{Εισαγωγή}
\label{ch:eisagogi}

\section{Κίνητρο}
\label{sec:kinitro}

Τα \en{Non-Fungible Tokens (NFTs)} είναι ψηφιακά περιουσιακά στοιχεία των οποίων η
μοναδικότητα και η κυριότητα πιστοποιούνται από ένα δημόσιο \en{blockchain}. Από την
τυποποίηση του \en{ERC-721} στο \en{Ethereum} \cite{erc721} έως τα σύγχρονα
αντικειμενοκεντρικά \en{blockchains} όπως το \en{Sui} \cite{sui-whitepaper}, τα \en{NFTs}
χρησιμοποιούνται για ψηφιακή τέχνη, συλλεκτικά αντικείμενα, εισιτήρια, αντικείμενα
παιχνιδιών και τίτλους ιδιοκτησίας. Η δευτερογενής αγορά τους, ωστόσο, οργανώθηκε εξαρχής
γύρω από ένα και μόνο πρότυπο συναλλαγής: την πώληση έναντι κρυπτονομίσματος. Πλατφόρμες
όπως το \en{OpenSea} \cite{opensea} και το \en{Blur} \cite{blur} στο \en{Ethereum}, ή το
\en{Tradeport} \cite{tradeport} και το \en{BlueMove} \cite{bluemove} στο \en{Sui}, προσφέρουν καταχωρήσεις σταθερής τιμής, δημοπρασίες και
προσφορές σε νόμισμα. Η μονάδα μέτρησης της αξίας είναι πάντοτε το \en{ETH} ή το
\en{SUI}.

Υπάρχει όμως μια ολόκληρη κατηγορία συναλλαγών που το πρότυπο αυτό δεν εξυπηρετεί: η
\emph{ανταλλαγή} (\en{barter}). Το πρόβλημα γίνεται φανερό με ένα απλό παράδειγμα. Έστω
ότι κάποιος έχει ένα μήλο και ένας φίλος του ένα πορτοκάλι, και ο καθένας θέλει το φρούτο
του άλλου. Η φυσική λύση είναι ο πρώτος να δώσει το μήλο και ο δεύτερος το πορτοκάλι. Σε
μια αγορά που γνωρίζει μόνο πωλήσεις, όμως, η ανταλλαγή αυτή δεν μπορεί καν να εκφραστεί:
πρέπει πρώτα ο καθένας να πουλήσει το φρούτο του έναντι του νομίσματος της χώρας και έπειτα
να αγοράσει από τον άλλον το φρούτο που τον ενδιαφέρει, με δύο ανεξάρτητες συναλλαγές,
διπλά τέλη και τον κίνδυνο να μεταβληθεί η τιμή στο ενδιάμεσο. Ακριβώς σε αυτή τη θέση
βρίσκεται σήμερα ένας συλλέκτης που κατέχει ένα αντικείμενο μιας συλλογής και επιθυμεί ένα
αντικείμενο διαφορετικής συλλογής: είτε πουλά και αγοράζει σε δύο βήματα, είτε
διαπραγματεύεται εκτός αλυσίδας μια απευθείας ανταλλαγή. Η δεύτερη επιλογή αναδεικνύει το
κλασικό πρόβλημα της δίκαιης ανταλλαγής: κάποιος από τους δύο πρέπει να στείλει πρώτος,
εμπιστευόμενος ότι ο άλλος θα ανταποκριθεί, ή πρέπει και οι δύο να εμπιστευθούν έναν
μεσάζοντα-θεματοφύλακα.

Το \en{Sui} προσφέρει τεχνικές προϋποθέσεις που καθιστούν το πρόβλημα αυτό ασυνήθιστα
ευχερές να λυθεί σωστά. Κάθε \en{NFT} είναι ένα αντικείμενο πρώτης τάξης με δικό του
αναγνωριστικό και ιδιοκτήτη, όχι μια εγγραφή σε έναν χάρτη μέσα σε ένα συμβόλαιο. Η γλώσσα
\en{Move} \cite{move-paper} διασφαλίζει στατικά ότι ένα τέτοιο αντικείμενο δεν μπορεί να
αντιγραφεί, να χαθεί ή να δημιουργηθεί από το πουθενά. Τα \en{Programmable Transaction
Blocks} επιτρέπουν σε μία συναλλαγή να συνθέσει πολλαπλές κλήσεις, και τα \emph{κοινόχρηστα
αντικείμενα} (\en{shared objects}) επιτρέπουν σε δύο άγνωστους μεταξύ τους χρήστες να
αλληλεπιδράσουν με την ίδια δομή δεδομένων χωρίς κεντρικό διαχειριστή. Με άλλα λόγια, μια
προσφορά ανταλλαγής μπορεί να γίνει η ίδια ένα αντικείμενο που κρατά δεσμευμένα τα
προσφερόμενα στοιχεία, και η αποδοχή της μια ατομική συναλλαγή που τα μεταβιβάζει.

\section{Αντικείμενο και στόχος}
\label{sec:stoxos}

Αντικείμενο της εργασίας είναι η σχεδίαση και η υλοποίηση ενός \en{NFT marketplace}
ανταλλαγών στο \en{Sui}, με κεντρικό χαρακτηριστικό τις \emph{προσφορές με \en{NFTs} από
διαφορετικές συλλογές}. Συγκεκριμένα, το σύστημα πρέπει να επιτρέπει σε έναν χρήστη:

\begin{itemize}[nosep]
  \item να υποβάλει προσφορά (\en{bid}) που αποτελείται από ένα ή περισσότερα \en{NFTs}
        οποιασδήποτε συλλογής και, προαιρετικά, από ποσό \en{SUI}, με αντάλλαγμα ένα
        συγκεκριμένο \en{NFT} που ανήκει σε άλλον χρήστη·
  \item να δεσμεύει τα προσφερόμενα στοιχεία κατά την υποβολή, ώστε ο αποδέκτης να γνωρίζει
        ότι η προσφορά είναι έγκυρη και εκτελέσιμη·
  \item ως κάτοχος του ζητούμενου \en{NFT}, να αποδεχθεί την προσφορά, οπότε η ανταλλαγή
        ολοκληρώνεται σε μία ατομική συναλλαγή, ή να την απορρίψει, οπότε τα δεσμευμένα
        στοιχεία επιστρέφουν στον προσφέροντα·
  \item ως προσφέρων, να ακυρώσει την προσφορά του οποιαδήποτε στιγμή πριν από την αποδοχή της.
\end{itemize}

Το σύστημα δεν διαθέτει διαχειριστή, θεματοφύλακα ή προνομιακά κλειδιά: η ορθότητα της
ανταλλαγής διασφαλίζεται αποκλειστικά από το έξυπνο συμβόλαιο και το μοντέλο εκτέλεσης του
\en{Sui}. Η υποστήριξη διαφορετικών συλλογών δεν επιτυγχάνεται με λίστες επιτρεπόμενων
τύπων ή με μετατροπή σε κάποιο κοινό ενδιάμεσο \en{token}, αλλά με γενικούς τύπους
(\en{generics}) στο επίπεδο του συμβολαίου: μια προσφορά έχει τύπο
\code{Bid<Offered, Requested>}, όπου οι δύο παράμετροι μπορούν να είναι οποιοιδήποτε
τύποι αντικειμένων με τις κατάλληλες ικανότητες. Έτσι ο ίδιος κώδικας εξυπηρετεί
ανταλλαγές εντός της ίδιας συλλογής, μεταξύ συλλογών, ή ακόμη και μεταξύ εντελώς
διαφορετικών εφαρμογών, με τον έλεγχο τύπων να γίνεται από την ίδια την εικονική μηχανή.

\section{Μεθοδολογία}
\label{sec:methodologia}

Η εργασία ακολούθησε τέσσερις φάσεις, αντίστοιχες με εκείνες που περιγράφουν οι
\cite{papias2023} και \cite{hadjichristofi2022}. Στην πρώτη φάση μελετήθηκαν το μοντέλο
αντικειμένων του \en{Sui} και η γλώσσα \en{Move}, ώστε να εντοπισθούν οι δυνατότητες που
απαντούν στο πρόβλημα που περιγράφηκε παραπάνω.
Στη δεύτερη φάση προσδιορίσθηκαν οι απαιτήσεις και σχεδιάσθηκε το πρωτόκολλο: οι ρόλοι, ο
κύκλος ζωής της προσφοράς, οι αναλλοίωτες συνθήκες ασφάλειας και η δομή των δεδομένων στην
αλυσίδα. Στην τρίτη φάση υλοποιήθηκε το έξυπνο συμβόλαιο με ανάπτυξη καθοδηγούμενη από
ελέγχους, και στη συνέχεια η \en{dApp}· και τα δύο αναθεωρήθηκαν αργότερα με βάση
δημοσιευμένα πρότυπα κώδικα και σχεδίασης διεπαφών. Στην τέταρτη φάση το σύστημα
αναπτύχθηκε σε τοπικό δίκτυο (και, σε πρώιμη έκδοση, σε \en{devnet}), επαληθεύθηκε από
άκρη σε άκρη και αξιολογήθηκε ως προς το κόστος των συναλλαγών και την αντοχή του.

\section{Δομή της εργασίας}
\label{sec:domi}

Η εργασία οργανώνεται ως εξής.

\begin{itemize}[nosep]
  \item Το Κεφάλαιο~\ref{ch:theoria} παρουσιάζει το θεωρητικό υπόβαθρο: τις βασικές
        έννοιες των \en{blockchains}, το αντικειμενοκεντρικό μοντέλο του \en{Sui}, τη
        γλώσσα \en{Move} και τα \en{NFTs}.
  \item Το Κεφάλαιο~\ref{ch:apaitiseis} καταγράφει τις απαιτήσεις και τη σχεδίαση του
        συστήματος, με περιπτώσεις χρήσης, διαγράμματα \en{UML} και τη μηχανή καταστάσεων
        της προσφοράς.
  \item Το Κεφάλαιο~\ref{ch:contract} περιγράφει την υλοποίηση του έξυπνου συμβολαίου σε
        \en{Move}, την ανάλυση ασφάλειας και τους ελέγχους.
  \item Το Κεφάλαιο~\ref{ch:dapp} περιγράφει την αρχιτεκτονική και την υλοποίηση της
        \en{dApp}.
  \item Το Κεφάλαιο~\ref{ch:epideixi} επιδεικνύει τη λειτουργία του συστήματος μέσα από
        τις οθόνες της εφαρμογής.
  \item Το Κεφάλαιο~\ref{ch:axiologisi} αξιολογεί το σύστημα ως προς το κόστος, την
        αντοχή και την κάλυψη των απαιτήσεων.
  \item Το Κεφάλαιο~\ref{ch:symperasmata} συνοψίζει τα συμπεράσματα και προτείνει
        μελλοντικές επεκτάσεις.
  \item Στο παράρτημα παρατίθεται ο πλήρης πηγαίος κώδικας των δύο \en{modules} του
        έξυπνου συμβολαίου και του σκηνικού των ελέγχων του.
\end{itemize}
